% TODO rimuovere
%\section{Implementazione del Sistema}


% OK DA FILIPPO
\begin{frame}{Assurance Engine ed Elasticsearch}
  \begin{itemize}


    \item L'Assurance Engine \alert{implementa} il \textbf{processo di assurance} precedente \alert{interagendo} con svariati \alert{sistemi di monitoring} (Elasticsearch,Prometheus)
    % prima erano astratti, questa è la concretizzazione
   
    \item \alert{API REST} esposte tramite OpenAPI per interagire con l'engine
    
    \item Integrazione di \alert{Elasticsearch} con l'Assurance Engine come \alert{strumento di monitoring}
    % per la raccolta, analisi e ricerca dei dati

    \begin{itemize}
      \item Efficienza nella \alert{raccolta continua} di dati, nella \alert{ricerca} e nell'\alert{analisi}
      
      %\item \alert{Scalabilità} del sistema
    \end{itemize}
  \end{itemize}
\end{frame}


\begin{frame}{Necessità e Vantaggi del Framework Intermediario}
  \begin{itemize}
    \item \alert{Estende} le capacità di Elasticsearch, utilizzando il linguaggio di programmazione \alert{Rust} 
    
    \item Infrastruttura \alert{request-reflector} 
    
    \item Fornisce \alert{scalabilità} e \alert{replicabilità senza coordinazione}:
    %L’Assurance Engine viene eseguito su un cluster Kubernetes
    
    \begin{itemize}
      \item Sistema deployato come \alert{una o più} istanze
      %\alert{Network load balancer} gestisce il traffico

      \item Istanze \alert{intercambiabili} e \alert{stateless} non necessitano di coordinazione aggiuntiva
      
      %\item Richieste gestibili da qualsiasi istanza, assicurando \alert{consistenza} e \alert{disponibilità}
    \end{itemize}
  \end{itemize}
\end{frame}


% ANCHE QUESTO SECONDO FILIPPO SI POTREBBE TAGLIARE
\begin{frame}{Funzionamento dell'Assurance Engine}

  \begin{enumerate}
    \item \alert{Richiesta} POST all'endpoint dell'Assurance Engine
    
    \item Valutazione di un contratto in \alert{tre fasi}:
    
    \begin{enumerate}
      \item \textbf{Measurement collection}: interrogazione dei \alert{servizi di monitoring} %tramite metriche
      
      \item \textbf{Contract evaluation}: valutazione dei \alert{contratti} 
      %tramite un confronto tra i parametri e le \alert{evidence} raccolte
      
      \item \textbf{Report composition}: restituzione del \alert{risultato} della valutazione %in formato JSON
    \end{enumerate}
  \end{enumerate}
\end{frame}
